% !TeX program = xelatex
% !TeX encoding = UTF-8
% =====================================================================
%  编译方式: XeLaTeX
%  说明: 全文已脱敏, 不含真实姓名 / 单位 / 联系方式 / 密钥
% =====================================================================
\documentclass[11pt,a4paper,openany]{ctexrep}

% ----------------------------- 基础宏包 -------------------------------
\usepackage[margin=2.5cm,top=2.8cm,bottom=2.8cm]{geometry}
\usepackage{fontspec}
\usepackage[table]{xcolor}
\usepackage{graphicx}
\usepackage{array}
\usepackage{booktabs}
\usepackage{longtable}
\usepackage{enumitem}
\usepackage{tikz}
\usepackage[most]{tcolorbox}
\usepackage{listings}
\usepackage{fancyhdr}
\usepackage{titlesec}
\usepackage{hyperref}
\usepackage{caption}

% ----------------------------- TikZ 库 --------------------------------
\usetikzlibrary{
  positioning, arrows.meta, calc, fit, backgrounds,
  shapes.geometric, shapes.misc, shadows, trees, matrix
}

% ----------------------------- 正式配色方案 -------------------------------
% 主色调：深蓝+深灰，符合正式报告规范
\definecolor{official_blue}{RGB}{0, 51, 102}
\definecolor{official_blue_light}{RGB}{217, 226, 243}
\definecolor{official_gray}{RGB}{51, 51, 51}
\definecolor{official_gray_light}{RGB}{242, 242, 242}
\definecolor{official_red}{RGB}{153, 0, 0}
\definecolor{official_green}{RGB}{0, 102, 51}
\definecolor{table_head}{RGB}{28, 56, 89}
\definecolor{table_alt}{RGB}{247, 248, 250}
\definecolor{code_bg}{RGB}{31, 31, 31}

% ----------------------------- 字体设置 -------------------------------
% 正式文档：正文宋体，标题黑体，英文 Times New Roman
\setmainfont{Times New Roman}
\setsansfont{Arial}
\setmonofont{Courier New}
\ctexset{
  chapter/format = \huge\heiti\bfseries\color{official_blue},
  section/format = \Large\heiti\bfseries\color{official_gray},
  subsection/format = \large\heiti\bfseries\color{official_gray}
}

% ----------------------------- 超链接 ---------------------------------
\hypersetup{
  colorlinks=true,
  linkcolor=official_blue,
  urlcolor=official_blue,
  citecolor=official_blue,
  pdftitle={智能在线银行系统-技术报告},
  pdfauthor={Dev Team}
}

% ----------------------------- 页眉页脚 -------------------------------
\pagestyle{fancy}
\fancyhf{}
\setlength{\headheight}{15pt}
\renewcommand{\headrulewidth}{0.6pt}
\renewcommand{\headrule}{\color{official_blue}\hrule height \headrulewidth}
\fancyhead[L]{\small\heiti\color{official_gray} 智能在线银行系统 · 技术报告}
\fancyhead[R]{\small\color{official_gray} \leftmark}
\fancyfoot[C]{\thepage}
\fancypagestyle{plain}{\pagestyle{fancy}}

% --------------------------- 章节标题样式 -----------------------------
\titleformat{\chapter}[hang]
  {\huge\heiti\bfseries\color{official_blue}}
  {\thechapter}
  {0.8em}
  {}
\titlespacing*{\chapter}{0pt}{0pt}{20pt}

\titleformat{\section}[hang]
  {\Large\heiti\bfseries\color{official_gray}}
  {\thesection}
  {0.6em}
  {}
\titlespacing*{\section}{0pt}{18pt}{10pt}

\titleformat{\subsection}[hang]
  {\large\heiti\bfseries\color{official_gray}}
  {\thesubsection}
  {0.5em}
  {}
\titlespacing*{\subsection}{0pt}{12pt}{8pt}

% --------------------------- tcolorbox 正式样式 ---------------------------
% 信息框
\newtcolorbox{infobox}[1][信息说明]{
  enhanced, breakable,
  colback=official_blue_light!30, colframe=official_blue,
  boxrule=0.8pt, arc=2pt, left=12pt, right=12pt, top=8pt, bottom=8pt,
  title={\textbf{#1}}, fonttitle=\heiti\bfseries\color{white},
  coltitle=white, attach boxed title to top left={xshift=10pt,yshift=-3mm},
  interior style={left color=official_blue_light!20, right color=white}
}

% 警告框
\newtcolorbox{warnbox}{
  enhanced, breakable,
  colback=red!5, colframe=official_red,
  boxrule=0.8pt, arc=2pt, left=12pt, right=12pt, top=8pt, bottom=8pt,
  interior style={left color=red!8, right color=white}
}

% 成果框
\newtcolorbox{achbox}{
  enhanced, breakable,
  colback=green!5, colframe=official_green,
  boxrule=0.8pt, arc=2pt, left=12pt, right=12pt, top=8pt, bottom=8pt,
  interior style={left color=green!8, right color=white}
}

% 代码框
\newtcolorbox{codebox}[1][终端命令]{
  enhanced,
  colback=white,               % 背景：亮色白底
  colframe=official_gray!80,  % 边框浅灰
  boxrule=1pt, arc=2pt, 
  left=8pt, right=8pt, top=22pt, bottom=8pt,
  title={\texttt{\color{gray!80}#1}}, 
  fonttitle=\bfseries,
  coltitle=gray!80, 
  attach boxed title to top left={xshift=10pt,yshift=-2.5mm},
  listing only,
  listing options={
    basicstyle=\ttfamily\small\color{black}, % 基础文字黑色
    keywordstyle=\color{blue!80},            % 关键字蓝色（linux亮色终端风格）
    commentstyle=\color{green!60},           % 注释绿色
    stringstyle=\color{red!80},              % 字符串红色
    showstringspaces=false, 
    language=bash, 
    frame=none
  }
}


% --------------------------- 列表样式 ---------------------------------
\setlist[itemize]{leftmargin=1.8em, itemsep=4pt, topsep=4pt}
\setlist[enumerate]{leftmargin=2.0em, itemsep=5pt, topsep=4pt}
\newcommand{\tagdot}[2]{\tikz[baseline=-0.5ex]\node[circle,fill=#1,inner sep=3pt]{};~\textbf{#2}}

% --------------------------- 代码清单样式 -----------------------------
\lstdefinestyle{yamlstyle}{
  basicstyle=\ttfamily\footnotesize\color{official_gray},
  backgroundcolor=\color{official_gray_light},
  keywordstyle=\color{official_blue}\bfseries,
  commentstyle=\color{official_green!70},
  frame=single, rulecolor=\color{official_gray!30},
  framesep=8pt, showstringspaces=false, breaklines=true, tabsize=2
}

% =====================================================================
%                               封 面
% =====================================================================
\begin{document}
\begin{titlepage}
  \centering
  \vspace*{1.5cm}
  
  % 抬头
  {\Large\heiti\color{official_gray} 课程设计技术报告}\\[0.8cm]
  {\color{official_blue}\rule{0.6\textwidth}{1.5pt}}\\[1.2cm]
  
  % 主标题
  {\fontsize{36pt}{42pt}\heiti\bfseries\color{official_blue} 智能在线银行系统}\\[0.6cm]
  {\fontsize{18pt}{22pt}\bfseries\color{official_gray} Smart Online Banking System}\\[0.8cm]
  
  % 副标题
  {\large\heiti\color{official_gray} AI驱动 · 前后端分离 · 全流程银行业务实践}\\[1.5cm]
  {\color{official_blue!60}\rule{0.4\textwidth}{1pt}}\\[1.5cm]
  
  % 项目信息表
  \begin{tcolorbox}[
    width=0.65\textwidth, center,
    colback=white, colframe=official_blue,
    boxrule=1pt, arc=3pt
  ]
    \centering
    \renewcommand{\arraystretch}{1.6}
    \begin{tabular}{>{\heiti\color{official_blue}}r l}
      项目类型 & 个人项目 · 前后端分离工程 \\
      团队规模 & 4 人开发组 \\
      担任角色 & 项目组长 \\
      交付规模 & 39 个 REST 接口 · 16 个页面 \\
      完成日期 & 2026 年 6 月 \\
    \end{tabular}
  \end{tcolorbox}
  
  \vfill
  
  % 底部
  {\small\color{official_gray!70} 本文档所有数据均为模拟数据 · 仅供学习与技术研究使用}
\end{titlepage}

% =====================================================================
%                               目 录
% =====================================================================
\tableofcontents
\thispagestyle{fancy}
\clearpage

% =====================================================================
%                          一、项目简介
% =====================================================================
\chapter{项目简介}

\begin{infobox}[一句话定位]
基于 \textbf{Spring Boot 3 + Vue 3 + 大模型} 的智能在线银行系统，实现银行业务全流程数字化与 AI 智能客服落地。
\end{infobox}

本项目为前后端分离工程，通过 \textbf{Spring AI} 接入大模型，提供智能客服、账户管理、转账交易、账单查询、贷款管理等功能。系统共包含 \textbf{39 个 REST 接口、16 个前端页面}，由 4 人小组协作完成；本人担任项目组长，负责架构设计、AI 客服模块开发与测试工作。

\begin{center}
\renewcommand{\arraystretch}{1.5}
\rowcolors{2}{table_alt}{white}
\begin{tabular}{>{\centering\arraybackslash}p{3cm} >{\centering\arraybackslash}p{3cm}
>{\centering\arraybackslash}p{3cm} >{\centering\arraybackslash}p{3cm}}
\rowcolor{table_head}\color{white}\textbf{后端类} & \color{white}\textbf{接口数} &
\color{white}\textbf{页面数} & \color{white}\textbf{团队规模}\\
\textbf{70+} & \textbf{39} & \textbf{16} & \textbf{4 人}\\
\end{tabular}
\end{center}

\begin{warnbox}
\textbf{特别说明：}本项目为课程设计实践，所有数据均为模拟数据，不涉及任何真实金融业务；系统中涉及的机构名称、业务数据与真实机构无关。
\end{warnbox}

% =====================================================================
%                          二、技术栈
% =====================================================================
\chapter{技术栈}

系统采用成熟稳定的主流技术体系，后端基于 Java 生态，前端基于 Vue 生态，AI 能力通过 Spring AI 独立封装。

\begin{center}
\renewcommand{\arraystretch}{1.42}
\rowcolors{2}{table_alt}{white}
\begin{longtable}{>{\raggedright\arraybackslash}p{3.2cm}
>{\raggedright\arraybackslash}p{6.2cm} >{\centering\arraybackslash}p{3.4cm}}
\rowcolor{table_head}\color{white}\textbf{分类} & \color{white}\textbf{技术名称} &
\color{white}\textbf{版本}\\
\endhead
开发语言 & Java & 17 LTS\\
后端框架 & Spring Boot & 3.2.5\\
安全框架 & Spring Security & 6.x\\
AI 框架 & Spring AI & 1.0.0-M4\\
ORM 框架 & MyBatis-Plus & 3.5.6\\
关系型数据库 & MySQL & 8.0\\
对象存储 & MinIO & 8.5.10\\
前端框架 & Vue & 3.4\\
UI 组件库 & Ant Design Vue & 4.2\\
状态管理 & Pinia & 2.1\\
前端构建 & Vite & 5.4\\
图表组件 & ECharts & 5.5\\
接口文档 & Knife4j & 4.4.0\\
项目构建 & Maven & 3.8+\\
测试框架 & JUnit & 5.x\\
\end{longtable}
\end{center}

% =====================================================================
%                          三、系统架构
% =====================================================================
\chapter{系统架构}

系统采用标准分层架构，前端通过 JWT 完成身份认证，后端按控制层、业务逻辑层、数据访问层分层设计，AI 能力以服务形式封装，与业务模块解耦。

\begin{center}
\begin{tikzpicture}[
  node distance=8mm and 14mm,
  box/.style={rectangle, rounded corners=3pt, minimum height=11mm,
    minimum width=58mm, align=center, font=\small\bfseries,
    text=white, drop shadow={shadow xshift=1.5pt, shadow yshift=-1.5pt, opacity=0.2}},
  data/.style={cylinder, shape border rotate=90, aspect=0.28,
    minimum height=14mm, minimum width=30mm, align=center,
    font=\small\bfseries, text=white},
  arr/.style={-{Stealth[length=3mm]}, line width=0.9pt, color=official_gray!60}
]
  \node[box, fill=official_green] (fe) {Vue 前端页面\\{\normalsize Vite · Ant Design Vue}};
  \node[box, fill=official_blue, below=of fe] (sec) {Spring Security\\{\normalsize JWT 认证 · RBAC 授权}};
  \node[box, fill=official_blue!80, below=of sec] (ctrl) {业务控制层 Controller};
  \node[box, fill=official_blue!60, below=of ctrl] (svc) {业务逻辑层 Service};
  \node[data, fill=official_blue!70, below left=10mm and 26mm of svc] (mysql) {MySQL\\数据库};
  \node[data, fill=official_gray!70, below=10mm of svc] (minio) {MinIO\\对象存储};
  \node[box, fill=official_red, minimum width=32mm, below right=10mm and 24mm of svc] (llm) {大模型服务};
  \draw[arr] (fe)--(sec);
  \draw[arr] (sec)--(ctrl);
  \draw[arr] (ctrl)--(svc);
  \draw[arr] (svc)--(mysql);
  \draw[arr] (svc)--(minio);
  \draw[arr] (svc)--(llm);
\end{tikzpicture}
\captionof{figure}{系统整体架构图}
\end{center}

\begin{infobox}[架构说明]
请求链路为：\textbf{前端 JWT 认证 → 控制层 → 逻辑层 → 数据层}；大模型调用独立封装，支持同步对话、SSE 流式输出与图片生成，异常时自动降级为本地兜底回复。
\end{infobox}

% =====================================================================
%                          四、功能模块
% =====================================================================
\chapter{功能模块}

系统围绕银行零售业务，划分为 8 个功能模块，实现业务全流程闭环。

\begin{center}
\renewcommand{\arraystretch}{1.45}
\rowcolors{2}{table_alt}{white}
\begin{longtable}{>{\centering\arraybackslash}p{1.2cm}
>{\raggedright\arraybackslash}p{3.4cm} >{\raggedright\arraybackslash}p{8.2cm}}
\rowcolor{table_head}\color{white}\textbf{序号} & \color{white}\textbf{功能模块} &
\color{white}\textbf{主要功能}\\
\endhead
1 & 用户管理 & 用户注册、登录认证、个人信息维护\\
2 & 账户管理 & 账户查询、充值、账户流水、余额管理\\
3 & 转账交易 & 行内转账、限额校验、交易记录查询\\
4 & 账单管理 & 月度账单自动生成、收支明细查询\\
5 & 贷款管理 & 贷款申请、审批、分期还款\\
6 & 支票管理 & 支票开具、兑付、作废（管理员）\\
7 & AI 智能客服 & 多会话管理、流式对话、上下文记忆、图片生成\\
8 & 管理后台 & 数据统计、用户管理、账户管理、贷款审批\\
\end{longtable}
\end{center}

% =====================================================================
%                          五、项目结构
% =====================================================================
\chapter{项目结构}

项目工程划分为后端工程、前端工程、数据库脚本与文档资源四大部分，整体结构如下（三层组织架构，直角连线）。

\begin{center}
\begin{figure}[htb]
	\centering
	\includegraphics[width=\linewidth]{figure1.png}
\end{figure}
\captionof{figure}{项目组织结构图}
\end{center}

% =====================================================================
%                          六、项目成果
% =====================================================================
\chapter{项目成果}

\begin{achbox}
\begin{enumerate}[label=\textbf{\color{official_green}\arabic*.}]
  \item 完成 \textbf{8 大功能模块}，实现银行业务全流程闭环，共交付 39 个接口、16 个页面；
  \item AI 客服支持\textbf{流式输出与多轮上下文记忆}，通过系统提示词约束、敏感词过滤、防诈骗提醒保障输出安全；
  \item 设计\textbf{大模型故障兜底机制}，服务异常时自动切换本地回复，保证系统稳定可用；
  \item 转账功能通过\textbf{数据库事务与限额校验}，保障资金数据一致性；
  \item 项目通过功能、接口及性能测试，\textbf{课程设计获评优秀}。
\end{enumerate}
\end{achbox}

\section{关键工程问题与解决}
\begin{itemize}
  \item \tagdot{official_blue}{大模型同步调用阻塞主流程}：采用响应式编程与 SSE 流式返回，降低首字等待时间；
  \item \tagdot{official_red}{输出越界与敏感信息诱导}：强约束 System Prompt + 敏感词过滤 + 防诈骗提醒三层防护；
  \item \tagdot{official_gray}{大模型服务不可用}：异常自动降级为本地关键词规则引擎，客服服务不中断；
  \item \tagdot{official_green}{转账并发数据一致性}：声明式事务 + 余额双重校验 + 限额前置拦截。
\end{itemize}

% =====================================================================
%                          七、快速启动
% =====================================================================
\chapter{快速启动}

\section{初始化数据库}
依次执行 \texttt{sql} 目录下的建表脚本与初始化数据脚本。

\section{配置参数（脱敏示例）}
\begin{lstlisting}[style=yamlstyle]
spring:
  datasource:
    url: jdbc:mysql://localhost:3306/online_bank?serverTimezone=Asia/Shanghai
    username: ${DB_USERNAME}
    password: ${DB_PASSWORD}
  ai:
    openai:
      api-key: ${LLM_API_KEY}
      base-url: https://ark.example.com/api/v3
jwt:
  secret: ${JWT_SECRET}
\end{lstlisting}

\begin{warnbox}
密钥、数据库口令等敏感信息\textbf{一律通过环境变量传入}，禁止明文写入配置文件或提交至代码仓库。
\end{warnbox}

\section{启动后端服务}
\begin{codebox}
mvn spring-boot:run
\end{codebox}

\section{启动前端服务}
\begin{codebox}
cd frontend
npm install
npm run dev
\end{codebox}

\begin{center}
\renewcommand{\arraystretch}{1.45}
\rowcolors{2}{table_alt}{white}
\begin{tabular}{>{\raggedright\arraybackslash}p{4cm} >{\raggedright\arraybackslash}p{8cm}}
\rowcolor{table_head}\color{white}\textbf{访问入口} & \color{white}\textbf{地址}\\
前端页面 & \url{http://localhost:3000}\\
接口文档 & \url{http://localhost:8080/doc.html}\\
健康检查 & \url{http://localhost:8080/actuator/health}\\
\end{tabular}
\end{center}

% =====================================================================
%                          八、免责声明
% =====================================================================
\chapter{免责声明}

\begin{infobox}[声明]
本项目仅用于学习交流与技术研究。系统中涉及的机构名称、业务数据均为模拟数据，与真实机构无关；因使用本项目内容所引发的任何后果，由使用者自行承担。
\end{infobox}

\begin{center}
\vspace{400pt}
{\color{official_blue}\rule{0.4\textwidth}{1pt}}\\[10pt]
{\small\color{official_gray} By BUAS}\\[4pt]
{\footnotesize\color{official_gray!70} Dev Team · 2026}
\end{center}

\end{document}
